\section{对今天 LLM 与 Agent 实践的启示}

\subsection{别把语言模型当成 ``万能语义机''}

这节课对今天工程实践者最重要的提醒之一，是不要因为模型输出流畅，就误以为语义问题已经被彻底解决。对实际系统来说，至少还要分清三件事：

\begin{enumerate}
    \item \textbf{表面流畅性}：模型能否说出自然、连贯、貌似合理的话
    \item \textbf{任务正确性}：模型是否真的完成了用户意图，而不是只生成一段好看的文本
    \item \textbf{语境适配性}：模型是否知道在这个组织、这个工具链、这个文化语境下该如何表达
\end{enumerate}

\begin{warningbox}{很多产品失败，不是因为模型太弱，而是因为把流畅性当成了正确性}
当系统从聊天扩展到搜索、代码、工作流或 agent 执行时，这种混淆会迅速放大。一个会写漂亮自然语言的模型，不等于一个会稳定完成任务的系统。
\end{warningbox}

\subsection{语言学视角为什么对 Agent 依然重要}

今天很多 agent 系统看起来已经走出传统 NLP 的边界，但它们仍然深受语言现象约束：

\begin{itemize}
    \item 用户指令常带有省略、歧义、语气和隐含前提
    \item 多轮协作要求系统理解篇章结构和对话状态，而不只是单轮句子
    \item 工具调用和最终回答之间，需要明确区分事实、假设、计划和执行结果
\end{itemize}

\begin{knowledgebox}{Agent 本质上把 NLP 的老问题转成了更高风险的问题}
过去一个指代错误可能只会让 QA 模型答错一句话；今天同样的指代错误，可能让 agent 打开错误文件、调用错误工具或者误解组织规则。语言理解问题没有消失，只是进入了代价更高的系统层。
\end{knowledgebox}

\subsection{给工程团队的评估清单}

如果把 Lecture 18 的思想转成实际系统评估，可以得到下面这张表：

\begin{table}[H]
\centering
\begin{tabular}{p{0.2\textwidth} p{0.28\textwidth} p{0.34\textwidth}}
\toprule
评估维度 & 不该只看的指标 & 更应该补充看的内容 \\
\midrule
语言能力 & benchmark 总分、主观流畅度 & 指代、语用、长上下文一致性、风格迁移 \\
泛化能力 & 单一公开数据集表现 & 领域迁移、任务改写、低资源语言与 OOD 测试 \\
系统可靠性 & 单次 demo 成功率 & 多轮失败模式、恢复能力、工具误调用分析 \\
社会影响 & 简单安全分类标签 & 偏见、版权、隐私、能耗与人工劳动成本 \\
\bottomrule
\end{tabular}
\caption{把 Lecture 18 的问题意识转成工程评估清单}
\end{table}

\subsection{为什么这门课到今天仍然不过时}

尽管 Lecture 18 讨论的是 NLP、语言学与哲学，看起来不像一门 ``how to build agents'' 的实战课，但它反而解释了为什么很多现代系统在强模型之上仍会失败：因为模型之外，还有表示、语境、目标、评测和社会后果这些更大的问题。

\begin{importantbox}{系统越强，越需要回到基本问题}
当模型能力弱时，我们忙着追求更高分；当模型能力开始足够强时，真正拉开差距的是谁更清楚地理解语言任务本身。Lecture 18 的价值就在这里：它强迫我们重新思考，究竟什么叫理解，什么叫可靠，什么叫可用。
\end{importantbox}

\subsection{本章小结}

对今天的 LLM 和 agent 实践而言，这节课最重要的价值不是提供一个新技巧，而是提供一套更高层的判断框架：把语言能力、任务成功、泛化边界和社会成本分开看，才能避免把强模型误当成完整系统。

%% ========== 结语 ========== 

